\documentclass[11pt]{article}
\usepackage{amsmath,amssymb}
\title{Rechecking what remains of an issued curve after recalibration}
\date{26 November 2026}
\begin{document}
\maketitle
\paragraph{What the adapter supplies.}
The pinned ROS-Industrial VDA adapter transforms the node request to
robot-local coordinates but passes the edge trajectory field through
unchanged. Its control-point type describes world coordinates. The
replay runs the unchanged selected dispatch, node and edge methods,
captures the callback, and carries the exact integer control points
from that callback into the separate corridor calculation. Three
calibrations check the distinct coordinate treatment. The source
does not show if a downstream planner uses the edge, when it converts
it, or which curve its controller executes. It is wrong to join node
and edge numbers as if both were already in one frame. The modeled
consumer below explicitly retains the curve after an initial
world-to-local conversion; its initial transform is the identity.

\paragraph{Question and interface.}
When can a controller keep the uncompleted suffix of the captured
local curve after a new calibration or a changed clearance envelope?
The model uses one rational quadratic B\'ezier span with positive
weights, representing a degree-two clamped three-control-point
trajectory. Let $p_i$ denote the retained coordinates; the field values
received through the callback are not silently substituted for $p_i$
when the initial calibration is nonidentity. A fixed revised map
$A u+b$ is augmented by an unknown offset $z$ in a given compact convex
polygon $K$. Each corridor face requires $n\mathbin\cdot x\le c$.
A lower bound $\ell$ on curve progress means that all actual future
reference parameters until this contract ends lie in $[\ell,1]$ and
that the consumed execution instance is the one supplying the curve.
The classifier is exact for that declared envelope: its universal
test ranges over every $t$ in the suffix and every $z$ in $K$.
Each failure witness attains a violating point in that envelope;
the model does not decide that the actual controller will visit it.
A certificate that begins beyond observed progress requires an extra
proof, such as a monotone progress witness and a measured cut.
A stale status flag, command receipt, reached goal, or the presence of a
path target is not that witness. The corridor must already include any
conservative margin for body footprint and tracking error and must stay
valid through the claimed horizon.

\paragraph{Exact reusable check.}
Write $D(t)=\sum_i w_i B_i(t)$ and
$N_n(t)=\sum_i w_i B_i(t)\,n\mathbin\cdot(Ap_i+b)$,
where $B_i$ are the quadratic Bernstein basis functions.
All $w_i>0$ and $\sum_i B_i(t)=1$ give $D>0$ on $[0,1]$.
Let $h_K(n)=\max_{z\in K}n\mathbin\cdot z$; its witness is a vertex,
including for a degenerate polygon. The exact face test is
\[
  (c-h_K(n))D(t)-N_n(t)\ \geq\ 0 \qquad(\ell\le t\le1).
\]
One tests the two endpoints and, only for a convex scalar quadratic
with its minimizer strictly inside the interval, that minimizer.
Thus no tolerance sampling is needed and exact tangency passes. The
accompanying Lean file checks denominator positivity, interval quadratic
nonnegativity, and the projected support equivalence; it does not
formalize the downstream interface, the VDA code, or Python execution.
The independent numerical-algebraic implementation in the replay
interpolates the quadratic from three points for each actual $K$ vertex
instead of using the support and control coefficients. All 1{,}200
test verdicts agree and every rejection has a concrete parameter,
face and offset witness.

With an unchanged span and weights, the local numerators and denominator
are unchanged. Updating a known affine transform uses the same captured
local coefficients and reevaluates affected rows. Changing only a
face bound or its projected uncertainty uses the old transformed
polynomial for that normal and adjusts the multiple of $D$.
Progress changes only the interval; increasing a justified bound does
not invalidate an earlier positive decision under identical physical
inputs. A new normal, uncertainty support, frame, or changed curve may
require a new validation; a passed face decision cannot be reused
across it without a comparison. Multiple frame corrections compose
with the capture-to-current coordinate map, not by adding coordinate
differences; each continued execution must satisfy the law of its cut.
This is semantic reuse and exact witnessed classification, not a
speedup over an equally informed direct planner.

\paragraph{Distinguishing examples.}
For the captured points $(0,0),(2,-4),(4,0)$ with unit weights, a
half-turn about $(4,0)$ keeps the endpoint and places both transformed
endpoints inside the modeled corridor. The transformed interior
has $y=8t(1-t)$ though: a ceiling of $1$ is violated at $t=1/2$.
The suffix is safe exactly for
$\ell\geq 1/2+\sqrt 2/4$. With a second change tightening the ceiling to
$1/2$, the previously safe suffix beginning at $0.86$ fails. The
paired unit-weight curve with middle control point $(2,2)$ reaches
exactly $y=1$ and is safe at the first ceiling, although its control
hull fails that test. The parameter lower bound and fixed calibration
are inputs to these examples, not outputs of the callback. They
separate wrong endpoint-only reuse from conservative retirement
based on an enclosing hull.

\paragraph{Omitted laws.}
These results do not establish motion, safe approach to the suffix,
dynamics, release of an old physical hold, or continuing viability
after another fault. The union of uncertain outcomes may involve
rotation or a curved/nonconvex safe set; the scalar quadratic test is
not asserted for those cases. No hazard has been induced in live
equipment. No stack or contract acceptance is claimed.
\end{document}
